\subsection*{用能量控制有限时间内的状态}
第二章保守系统中的机械能也可用于延拓。这里不要求势能控制位置趋于无穷的行为：势能有下界即可先控制速度，再控制有限时间内的位置。
\begin{theorem}[能量与全局延拓]\label{thm:energy-global-extension}
设 $U\in C^2(\R^n)$ 有下界，$m_i>0$。则
$m_i\ddot x_i=-\partial_iU(x)$ 的每个解都能双向延拓到整条时间轴。
\end{theorem}
\begin{proof}
把二阶方程写成状态变量 $(x,v)$ 的一阶方程
$\dot x_i=v_i,\ \dot v_i=-\partial_iU(x)/m_i$。其向量场为 $C^1$，
所以每个初值有唯一最大解。沿该解计算
\[
 \frac{d}{dt}\left(\frac12\sum_i m_iv_i^2+U(x)\right)
 =\sum_i v_i\bigl(m_i\dot v_i+\partial_iU(x)\bigr)=0.
\]
设常能量为 $E_0$，取 $C$ 使 $U(x)\geq-C$，再令 $m_*:=\min_i m_i>0$。
于是
\[
 \begin{aligned}
 \frac{m_*}{2}|v(t)|^2&\leq\frac12\sum_i m_iv_i(t)^2\\
 &=E_0-U(x(t))\\
 &\leq E_0+C.
\end{aligned}
\]
因此存在与时间无关的 $M$，使 $|v(t)|\leq M$。
若最大正向存在端点 $\beta$ 有限，则对 $t_0\leq t<\beta$，
\[
 |x(t)|\leq |x(t_0)|+\int_{t_0}^t|v(s)|\,ds
 \leq |x(t_0)|+M(\beta-t_0).
\]
所以 $(x(t),v(t))$ 始终留在整个状态空间中的一个紧球内。
命题~\ref{prop:compact-extension-vector} 说明解可越过 $\beta$ 延拓，与最大性矛盾。
负向时间同样有
$|x(t)|\leq|x(t_0)|+M|t-t_0|$，故负向端点也不可能有限。
两个方向的最大端点分别为 $+\infty$ 和 $-\infty$。
\end{proof}
